%% -*- coding: utf-8 -*-


\section{形容词、副词和数词}

\subsection{形容词位置：前置或后置}

\subsection{形容词顺序}

\begin{center}
  \begin{talltblr}[
    label = none, % Removes the "Table X:" text
    ]{
    vlines, hlines,
    row{1}={font=\large\bfseries, bg=azure9},
    rows={halign=c, valign=m},
    columns={halign=c, valign=m}
    }
    \SetCell[c=9]{c}形容词的排序规则                                                             \\
    数量限定词 & 观点      & 大小  & 形状  & 新旧 & 颜色 & 国籍或产地 & 材料      & 被修饰的名词 \\
    two        & beautiful & small & round & new  & blue & Chinese    & porcelain & vases        \\
    \SetCell[c=9]{c}两只小巧精美、崭新的、蓝色的中式圆瓷瓶\\
  \end{talltblr}
\end{center}

\subsection{\(-ed\)形容词和\(-ing\)形容词}

\subsection{形容词比较级和最高级的形式}

\begin{enumerate}
  \item 规则变化
  \item 不规则变化
\end{enumerate}

\subsection{形容词比较级和最高级的用法}

\subsection{形容词表达感受和观点的句式}

\subsection{时间、地点、方式及频率副词}

\subsection{程度、疑问、关系、连接及其他副词}

\subsection{副词修饰整句话}

\subsection{副词的比较级和最高级}

\subsection{比较级的强调用法}

\subsection{\(as\)比较结构}

\subsection{常用比较级的句式}

\subsection{描述相似}

\subsection{比较数量}

\subsection{“too \(\ldots\) to” 结构 }

\subsection{数词倍数表达}

\subsection{数词 \(double\)}

\subsection{描述数量：可数、不可数名词}

\subsection{\(few\)和\(little\)}




%%% Local Variables:
%%% TeX-engine: xetex
%%% TeX-master: "../IELTS_300_Sentences"
%%% End:

% LocalWords:  vlines hlines halign valign xetex LocalWords bg
